% Teil II: der durchgerechnete Weg. Reihenfolge beibehalten.

\subsection*{Gegeben und gesucht}

Gegeben: <Größen mit Einheiten, benutztes Bezugssystem>.
Gesucht: <Größe mit Einheit>.

\subsection*{Ansatz}

<Ein Satz Begründung: welches Gesetz, welche Symmetrie, welche Erhaltungsgröße,
und warum dieser Weg der kürzeste ist.>

\subsection*{Rechnung}

<Symbolisch bis zum Schluss, Zahlen erst im letzten Schritt.>

\begin{equation}
  <Zusammenhang>
\end{equation}

\begin{ergebnis}
\begin{equation*}
  <Endergebnis mit Einheit>
\end{equation*}
\end{ergebnis}

\begin{kontrolle}
\textbf{Einheiten:} <Einheitenrechnung, die zum gesuchten Größe passt>.
\textbf{Grenzfall oder Querprobe:} <zweiter Weg oder Extremfall>.
\textbf{Größenordnung:} <ist die Zahl plausibel>.
\end{kontrolle}

\begin{merksatz}
<Was man aus dieser Aufgabe auf andere übertragen kann, in ein bis zwei Sätzen.>
\end{merksatz}
